\newpage
\section*{Aufgabenstellung}
\addcontentsline{toc}{section}{Aufgabenstellung}
\textbf{Titel:} Potenziale und Herausforderungen Digitaler Zwillinge in der Fabrikplanung mit NVIDIA Omniverse\\[0.5em]
\textbf{Bearbeiter:} Tobias Küster\\
\textbf{Erstprüfer:} Prof. Dr.-Ing. P.\ Diersen, Labor für Produktentwicklung und Digitale Fabrik\\
\textbf{Zweitprüfer:} Prof. Dr.-Ing. A.\ Vendl\\
\textbf{Betreuer LFP:} A.\ Ernst\\
\textbf{Ausgabe:} 19.06.2026\\[1em]

Industrie 4.0 und cyber-physische Produktionssysteme stellen die Fabrikplanung vor neue Anforderungen. Hohe Variantenvielfalt, kurze Produktzyklen und steigende Wandlungsfähigkeit verlangen dynamische statt statischer Planungsansätze. Digitale Zwillinge fungieren dabei als zentrale Technologie durch ihr kontinuierliches Abbild realer Systeme. Mit NVIDIA Omniverse und Isaac Sim entstehen Plattformen, die durch offene Standards, physikalische Genauigkeit und fotorealistisches Rendering neue Maßstäbe in der Simulation setzen. Daher sollen im Labor für Produktentwicklung und Digitale Fabrik (LFP) der Hochschule Hannover von zwei realen Anlagen Digitale Zwillinge erstellt werden.

\textbf{Teilziele:}
\begin{itemize}
    \item Recherche notwendiger theoretischer Grundlagen (Digitalisierungsstufen, Isaac Sim, Middleware OPC UA / MQTT / ROS\,2).
    \item Entwicklung Digitaler Zwillinge:
    \begin{enumerate}
        \item Rundtakttisch \enquote{Maze Runner} (Fertigung eines Produktes)
        \item Gelenkarmroboter \enquote{Dobot Magician} (Pick and Place Szenario)
    \end{enumerate}
    \item Inbetriebnahme und Optimierung des Maze Runner und Dobot Magician
    \item Reproduzierbarkeitsdokumentation (PowerPoint, Video, VR/AR).
    \item Ableitung von Potenzialen, Herausforderungen und Handlungsempfehlungen.
    \item Vollständige Dokumentation (max.\ 65 Seiten) und digitale Übergabe an das LFP; Open-Access-Veröffentlichung über die Bibliothek der HsH angedacht.
\end{itemize}
